% ===== 原文 pp.46--47（PDF pp.47--48）=====
% [原文 p.46]
且 $q^{\frac13}\Lambda\ll\mu\ll\Lambda$ 时处于相 II（对应 $\CP^2$）。让我们看看镜像描述里的情形是什么样子。把早先的计算翻译过来，我们看到有三个临界点满足 $x_1\sim x_2$、$x_1^3\sim q(\Lambda/\mu)^3$，以及一个单独的临界点满足 $x_1\sim-\Lambda/\mu$、$x_2\sim q\Lambda/\mu$。对于所考虑的标度窗口，我们看到前三个临界点处 $|x_1|\ll1$、$|x_2|\ll1$、$|q_1q_2|/|x_1x_2|\ll1$，而最后那个单临界点处 $|x_1|\gg1$、$|x_2|\ll1$、$|q_1q_2|/|x_1x_2|\gg1$。也就是说，前三个点对应的 $(\Re(y_1),\Re(y_3))$ 落在 $\CP^2$ 的矩多面体内（$\zeta=\zeta^1+\zeta^2$），而最后那个单点落在其外（见图 6）。这就清楚地表明，前三个临界点对应模型 $\CP^2$ 部分的真空。若取 $q\to0$ 并保持 $\Lambda_{\mathbb{P}^2}^3=q\Lambda^3$ 固定，把纯 $\CP^2$ 从 $\mathrm{H}_1$ 中挑出来，则超势变为

\begin{equation}
	\tW = \mu\left(x_1+\frac{(\Lambda_{\mathbb{P}^2}/\mu)^3}{x_1x_2}+x_2+\frac{q^{\frac13}(\Lambda_{\mathbb{P}^2}/\mu)^2}{x_2}\right)\longrightarrow \mu\left(x_1+\frac{(\Lambda_{\mathbb{P}^2}/\mu)^3}{x_1x_2}+x_2\right). \tag{5.37}
\end{equation}

这确实正是 $\CP^2$ 的镜像超势。

\para{$n=2$}

模型 $\mathrm{H}_2$ 在 $q_1\ll1$ 时被等同于靶空间为 $\mathbb{F}_2$ 的 sigma 模型，在 $q_1\gg1$ 时则被等同于靶空间为 $\WP^2_{1,1,n}$ 的 sigma 模型。取 $q_2=(\Lambda/\mu)^2$，镜像超势写为

\begin{equation}
	\tW = \mu\left(x_1+\frac{q_1(\Lambda/\mu)^4}{x_1x_2^2}+x_2+\frac{(\Lambda/\mu)^2}{x_2}\right). \tag{5.38}
\end{equation}

临界点为

\begin{equation}
	x_1 = \frac{\epsilon_2\epsilon_1q_1^{\frac12}\Lambda/\mu}{(1+2\epsilon_1q_1^{\frac12})^{\frac12}},\qquad x_2 = \epsilon_2(1+2\epsilon_1q_1^{\frac12})^{\frac12}\Lambda/\mu,\qquad \epsilon_1,\epsilon_2\in\{\pm1\}. \tag{5.39}
\end{equation}

在高能 $\mu\gg\Lambda$ 且 $q_1\ll1$ 时，它们满足 $|x_1|\ll1$、$|x_2|\ll1$、$|q_1q_2^2|/|x_1x_2^2|\ll1$、$|q_2|/|x_2|\ll1$，也就是说 $(\Re(y_1),\Re(y_3))$ 落在 $\mathbb{F}_2$ 的矩多面体内（图 8-右）。在高能 $\mu\gg\Lambda$ 且 $q_1\gg1$ 时，它们满足 $|x_1|\ll1$、$|x_2|\ll1$、$|q_1q_2^2|/|x_1x_2^2|\ll1$，也就是说 $(\Re(y_1),\Re(y_3))$ 落在 $\WP^2_{1,1,2}$ 的矩多面体内（图 8-左）。事实上，这是我们已经见过的一般事实的一个推论——在高能下，只要临界点在重整化变量中位于有限处，$(\Re(y_1),\Re(y_3))$ 就遍历 UV 靶空间的矩多面体。

% [Tutor Note: 由 (5.38) 解 (5.39) 可完整复核：令 $A=q_1(\Lambda/\mu)^4$、$B=(\Lambda/\mu)^2$，则 $\partial_{x_1}\tW=0$ 给出 $x_1^2x_2^2=A$，$\partial_{x_2}\tW=0$ 给出 $-2A/(x_1x_2^3)+1-B/x_2^2=0$；两式联立即 $x_2^2=B+2\sqrt{A}=(1+2\epsilon_1q_1^{\frac12})(\Lambda/\mu)^2$，其中 $\epsilon_1$ 标记 $\sqrt{q_1}$ 的两支，$\epsilon_2$ 标记 $\sqrt{x_2^2}$ 的两支；把 $x_1=\sqrt{A}/x_2$ 代回并用 $1/\epsilon_2=\epsilon_2$ 即得 (5.39) 的分子 $\epsilon_2\epsilon_1q_1^{\frac12}$。另：原文 §5 通篇写 "momentum polytope"，与 §4.5 的 "moment polytope"（矩多面体）是同一个对象，译文统一作"矩多面体"。]

\para{$n\ge3$}

模型 $\mathrm{H}_n$ 被等同于靶空间为 $\WP^2_{1,1,n}$ 的 sigma 模型。我们回忆，对于具有 $q_1=(\Lambda/\mu)^{2-n}$ 与 $q_2=q(\Lambda/\mu)^2$ 且 $q\ll1$ 的理论，在 $q^{\frac12}\Lambda\ll\mu\ll\Lambda$ 时处于相 I（对应 $\mathbb{F}_n$）。让我们看看镜像描述里的情形是什么样子。把早先的计算翻译过来，我们看到有四个临界点满足 $x_1\sim\epsilon_1q^{\frac{n}{4}}\Lambda/\mu$、$x_2\sim\epsilon_2q^{\frac12}\Lambda/\mu$，以及 $(n-2)$ 个临界点满足 $x_1^{n-2}\sim(-n)^{-n}(\Lambda/\mu)^{n-2}$、$x_2\sim-q(\Lambda/\mu)^2/x_1$。对于所考虑的标度窗口，我们看到前四个临界点处 $|x_1|\ll1$、

% [原文 p.47]

$|x_2|\ll1$、$|q_1q_2^n|/|x_1x_2^n|\ll1$、$|q_2|/|x_2|\ll1$，而后 $(n-2)$ 个临界点处 $|x_1|\gg1$、$|x_2|\ll1$、$|q_1q_2^n|/|x_1x_2^n|\gg1$、$|q_2|/|x_2|\gg1$。也就是说，前四个点对应的 $(\Re(y_1),\Re(y_3))$ 落在 $\mathbb{F}_n$ 的矩多面体内，而后 $(n-2)$ 个点落在其外（见图 8-右）。这就确定地表明，前四个临界点对应模型 $\mathbb{F}_n$ 部分的真空。然而，正如我们先前讨论过的，并不存在把纯 $\mathbb{F}_n$ 从 $\mathrm{H}_n$ 中分离出来的极限。相反，让我们看看在 $q\to0$ 且保持 $\Lambda$ 固定的极限下超势会发生什么，这个极限把 $|\phi_3|^2\sim\zeta^2$ 的理论约化为 $\mathrm{T}^{U(1)}_{2,n}(0)$：

\begin{align}
	\tW &= \mu\left(x_1+\frac{q_1q_2^n}{x_1x_2^n}+x_2+\frac{q_2}{x_2}\right)
	= \mu\left(x_1+\frac{(\Lambda/\mu)^{2-n}}{x_1\tx_2^n}+q(\Lambda/\mu)^2\tx_2+\frac{1}{\tx_2}\right)\nonumber\\
	&\longrightarrow \mu\left(x_1+\frac{(\Lambda/\mu)^{2-n}}{x_1\tx_2^n}+\frac{1}{\tx_2}\right). \tag{5.40}
\end{align}

这里我们使用变量 $x_1$ 与 $\tx_2:=x_2/q_2$，因为我们看的是场空间中 $|\phi_3|^2\sim\zeta^2$ 的区域。这恰恰就是模型 $\mathrm{T}^{U(1)}_{2,n}(0)$ 的镜像超势。

\para{模型 $\mathrm{T}^{U(1)}_{N,d}$}

作为检验，让我们求 (5.34) 的临界点。当 $d=N$（CY 情形）时，原点处有唯一的临界点，除非 $\ee^t=(-N)^N$——那是判别轨迹，见 (4.33)。当 $d\neq N$ 时，原点是孤立临界点，但原点之外还有 $|N-d|$ 个临界点（按模掉轨形群之后计数）。这些别的临界点对应有质量的 Coulomb 真空。在原点处的临界点上，当 $d<N$（相应地 $d>N$）时，(5.34) 的最后一项与前 $N$ 项相比是更高阶（相应地更低阶）的，因而无关（相应地相关）。因此，在深红外（相应地紫外），原点处的理论是 LG 轨形

\begin{equation}
	\tW = \tx_1^d+\cdots+\tx_N^d/(\Z_d)^{N-1}. \tag{5.41}
\end{equation}

这就是原模型中相应理论的镜像，也就是说，LG 轨形 $W=f(x_1,\dots,x_N)/\Z_d$。的确，当 $f$ 是 Fermat 型 $f(x)=x_1^d+\cdots+x_N^d$ 时这是成立的，这就是众所周知的 Greene–Plesser 镜像对称（Berglund–Hübsch 的 LG-LG 镜像对称的一个例子……）。

% [Tutor Note: (5.41) 中的斜杠表示对群取轨形，而非除法；其群与 (5.34) 下方 "$(\C^N/(\Z_d)^{N-1},(5.34))$" 相同。阶数核对：要求 Fermat 项 $\tx_1^d+\cdots+\tx_N^d$ 不变只需 $\omega_i^{d}=1$（$i=1,\dots,N$），而要求乘积项 $\tx_1\cdots\tx_N$（即 (5.34) 的 $\ee^{t/d}\tx_1\cdots\tx_N$）不变还需 $\prod_i\omega_i=1$；后者把 $(\Z_d)^N$ 裁掉一个 $\Z_d$，故群阶为 $d^{N-1}$，正是 $(\Z_d)^{N-1}$。另：原文 p.47 末尾写作 "Berglund-Hubsch"，通行拼写为 Berglund–Hübsch（见术语表）。原文 p.46 的 "a single critical points"、"This make it certain"、"there is no limit that isolate" 均为英语语法笔误，译文按原意处理，未改动数学内容。]
